이중 연결 요소(정점-서로소 이중 연결 요소)는 그래프를 트리로 압축하는 중요한 도구 중 하나이다. 임의의 무향 연결 그래프는 이중 연결 요소라 불리는 것으로 묶을 수 있고, 각 이중 연결 요소를 정점으로 한 트리로 압축할 수 있다. 이 절에서는 단절점과 단절선, 그리고 이중 연결 요소를 다룬다.

\section{단절점}
단절점은 무향 그래프에서 한 정점과 연결된 간선들을 제거했을 때 그래프가 둘 이상으로 나뉘는 정점을 말한다. 이러한 정점들을 구하는 문제를 생각해보자.

\begin{example} \label{ex: articulation-point}
    연결 그래프 $G=(V, E)$가 주어질 때, $V$의 각 원소 $v$에 대해 다음 조건을 만족하는지 판단하라.

    조건: $E$의 원소들 중 $v$와 다른 한 정점을 연결하는 간선들을 모두 제거한 집합 $E^\prime$에 대해 $G^\prime = (V - \{ v \}, E^\prime)$이 연결 그래프가 아니다.
\end{example}

다음과 같은 관찰을 기반으로 하면 단절점은 DFS를 이용하여 선형 시간에 구할 수 있다.

\begin{obs}
    시작 정점이 아닌 정점 $v$에 대해 $v$를 방문하기 이전에 방문한 정점 $a$와 $v$를 방문한 이후 방문한 정점 $b$가 연결된 순서쌍 $(a, b)$가 존재하면 $v$는 단절점이 아니다. 즉, $v$가 단절점이라는 것은 DFS order에서 $v$ 앞에 있는 임의의 정점과 $v$ 뒤에 있는 임의의 정점을 골랐을 때 연결되지 않았다는 것과 필요충분관계이다.
\end{obs}

즉, DFS order에서 자신의 자식 노드들이 방문한 정점(자신을 방문한 이후에 확인한 정점)들 중 자신보다 일찍 방문한 정점이 있다면 그 정점은 단절점이 아니고, 나머지 정점은 단절점이다.

단, 이러한 논의는 시작 정점에 대해서는 적용할 수 없는데, 시작 정점은 자신보다 일찍 방문한 정점이 없기 때문이다. 
하지만 시작 정점도 단절점이 될 수 있는데, 시작 정점에서 출발한 DFS가 종료되고 다시 시작 정점으로 돌아왔을 때 아직 방문하지 않은 정점이 남아있다면 둘 이상의 닫힌 부분이 있다는 뜻이므로 시작 정점이 단절점이 된다. 즉, 시작 정점의 경우 시작 정점에서 둘 이상의 재귀 함수 호출이 있으면 단절점이고, 그렇지 않으면 단절점이 아니다.

다음 \Cref{code: articulation-point}는 이 과정을 구현한 DFS 코드이다.

\begin{code} \label{code: articulation-point}
\Cref{ex: articulation-point}의 DFS 부분 코드
\begin{lstlisting}
int r = 0;
int solve(int v) {
    ord[v] = ++r;
    int fast = 1e9, child = 0;
    
    for (int i = 0; i < adj[v].size(); i++) {
        int u = adj[v][i];
        if (ch[u] == 1) fast = min(fast, ord[u]);
        else {
            child++;
            ch[u] = 1;
            int prev = solve(u);
            if (v != 1 && prev >= ord[v]) iscut[v] = 1;
            fast = min(fast, prev);
        }
    }
    
    if (v == 1 && child >= 2) iscut[v] = 1;
    
    return fast;
}
\end{lstlisting}
\end{code}

\section{단절선}
단절선은 단절점과 비슷한 개념이지만 정점이 아닌 간선을 제거하는 것이다. 즉, 무향 그래프에서 한 간선을 제거했을 때 그래프가 둘 이상으로 나뉘는 간선을 말한다. 이러한 간선들을 구하는 문제를 생각해보자.

\begin{example} \label{ex: articulation-bridge}
    연결 그래프 $G=(V, E)$가 주어질 때, $E$의 각 원소 $e$에 대해 다음 조건을 만족하는지 판단하라.

    조건: $G^\prime = (V, E - \{ e \})$이 연결 그래프가 아니다.
\end{example}

다음과 같은 관찰을 기반으로 단절선도 DFS를 이용하여 선형 시간에 구할 수 있다.

\begin{obs}
    어떤 간선 $e$에서 시작 정점에 가까운 정점을 $v$라고 하자. $v$로 도달하는 경로 중 DFS order에 따라 온 경로가 아닌 다른 경로에서 $v$보다 방문 순서가 이른 정점이 없으면 그 간선 $e$는 단절선이다.
\end{obs}

구현 자체는 단절점과 크게 다르지 않다. 간선의 시작 정점과 가까운 정점 $v$가 $v$를 방문한 이후에 확인한 정점들 중 \textbf{$v$의 부모가 아니고} 자신보다 일찍 방문한 정점이 있다면 그 간선은 단절선이 아니고, 나머지 간선은 단절선이다.

구현은 다음 \Cref{code: articulation-bridge}를 참고하라.

\begin{code} \label{code: articulation-bridge}
\Cref{ex: articulation-bridge}의 DFS 부분 코드
\begin{lstlisting}
int r = 0;
int solve(int v, int p) {
    ord[v] = ++r;
    int fast = 1e9;
    
    for (int i = 0; i < adj[v].size(); i++) {
        int u = adj[v][i].first, j = adj[v][i].second;
        if (u == p) continue;
        if (ch[u] == 1) fast = min(fast, ord[u]);
        else {
            ch[u] = 1;
            int prev = solve(u, v);
            if (prev > ord[v]) iscut[j] = 1;
            fast = min(fast, prev);
        }
    }
    
    return fast;
}
\end{lstlisting}
\end{code}

\section{이중 연결 요소}
이중 연결 요소는 단절점과 관련된 개념으로 한 정점과 연결된 간선을 제거했을 때 서로 연결되어 있는 가장 큰 부분 그래프를 말한다. 여기서 단절점을 기준으로 구분하면 각각의 부분 그래프가 이중 연결 요소가 되며 이러한 방식으로 그래프를 나눌 수 있다.

\begin{definition}[이중 연결 요소]
    연결 그래프 $G=(V, E)$의 이중 연결 요소는 $G$의 부분 연결 그래프이며 단절점이 없는 그래프 중 자신을 부분 그래프로 갖는 $G$의 이중 연결 요소가 없는 그래프를 말한다.
\end{definition}

한 그래프의 이중 연결 요소는 정점은 겹칠 수 있으나 간선은 겹치지 않는다. 따라서 구현에서 이중 연결 요소는 간선으로 저장한다.

구체적인 구현은 단절점의 구현에 간선을 묶어주는 부분만 추가하면 된다. 현재 보고 있는 점에 대해 연결된 간선이 아직 추가하지 않은 간선이면 스택에 간선을 넣는다. 지금 보고 있는 점이 단절점이라면 이 시점까지 확인한 부분이 하나의 이중 연결 요소이므로 스택에 넣은 간선 중 현재 보고 있는 간선까지 하나의 이중 연결 요소로 묶으면 된다.
구현 코드는 \Cref{code: bcc}를 참고하라.

\begin{code} \label{code: bcc}
이중 연결 요소의 DFS 부분 코드
\begin{lstlisting}
stack< pair<int, int> > stk;
vector< vector< pair<int, int> > > bcc;

int r = 0;
int solve(int v) {
    ord[v] = ++r;
    int fast = 1e9;
    
    for (int i = 0; i < adj[v].size(); i++) {
        int u = adj[v][i];
        if (ord[u] > ord[v]) stk.push({ord[v], ord[u]});
        if (ch[u] == 1) fast = min(fast, ord[u]);
        else {
            ch[u] = 1;
            int prev = solve(u);
            if (prev >= ord[v]) {
                vector< pair<int, int> > _bcc;
                while (stk.top() != {ord[v], ord[u]}) {
                    _bcc.push_back(stk.top());
                    stk.pop();
                }
                _bcc.push_back({ord[v], ord[u]});
                stk.pop();
                bcc.push_back(_bcc);
            }
            fast = min(fast, prev);
        }
    }
    
    return fast;
}
\end{lstlisting}
\end{code}

\section{블록-컷 트리}
이중 연결 요소를 바탕으로 그래프를 트리로 압축할 수 있다. 각 이중 연결 요소를 블록으로, 단절점을 컷으로 표현하면 블록과 컷을 연결하는 간선만으로 이루어진 블록-컷 트리를 만들 수 있다. 정확히는 각 단절점을 컷 정점으로, 각 이중 연결 요소를 블록 정점에 대응시키고 단절점이 이중 연결 요소의 정점 집합에 포함되는 컷과 블록을 연결하여 트리를 만든다.

블록-컷 트리를 사용하는 대표적인 문제를 하나 보자.
\begin{example} \label{ex: block-cut}
    어떤 무향 연결 그래프 $G=(V, E)$가 주어질 때, 다음 두 쿼리를 처리하는 프로그램을 작성하라.

    \begin{itemize}
        \item $u$와 $v$를 잇는 간선이 $E$에 포함되어 있을 때 $u$와 $v$를 잇는 간선을 제거한 후 $u$와 $v$가 연결되어 있는지 판단하라.
        \item $w$를 포함하는 모든 간선을 $E$에서 제거한 후 $u$와 $v$가 연결되어 있는지 판단하라.
    \end{itemize}

    (\url{https://www.acmicpc.net/problem/1734})
\end{example}

두 쿼리 모두 같은 방식으로 해결할 수 있다. 일단 쿼리 1에서는 $u$와 $v$를 잇는 간선이 단절선이 아니면 무조건 연결되어 있고 쿼리 2에서는 $w$가 단절점이 아니면 무조건 연결되어 있다.
이제 블록-컷 트리에서 $u$와 $v$를 잇는 유일한 경로를 생각하자. 이 경로 위에 $u$와 $v$를 잇는 간선을 포함하는 이중 연결 요소의 블록이 존재하면, 단절점 $w$에 해당하는 컷이 존재하면 $u$에서 $v$로 가는 경로가 절단되므로 더 이상 연결되지 않게 된다. 
이는 트리에서 최소 공통 조상을 사용하면 구현할 수 있으나, 아직 배우지 않았으므로 그 전까지를 구하는 것을 목표로 해보자. 이를 실제로 구현하는 문제는 \Cref{pb: block-cut}에 있다.

\section{연습문제 A}
\begin{problem}
    이중 연결 요소를 구하면 단절선도 구해낼 수 있다. 이중 연결 요소로부터 단절선을 구하는 방법을 생각해보고, 코드로 구현하라.
\end{problem}
\begin{problem} \label{pb: block-cut}
    \Cref{ex: block-cut}을 해결하는 코드를 작성하라. 배우지 않은 부분은 생략하여도 좋다.
\end{problem}
\begin{problem}
    \Cref{code: bcc}에서 각 이중 연결 요소에 포함되는 정점도 저장하도록 코드를 수정하여라.
\end{problem}

\section{연습문제 B}

\section{연습문제 C}